\documentclass[../../main]{subfiles}

\begin{document}

\section{Imagem Direta e Imagem Inversa}
\label{sec:matematica:funcoes:imagem-direta-e-inversa}

\subsection{Imagem Direta}

\begin{defn}[Imagem Direta]
	\label{def:matematica:funcoes:imagem-direta-e-inversa:imagem-direta}

	Seja

	\[
		f:A\to B
	\]

	e seja

	\[
		X\subseteq A.
	\]

	A imagem direta de \(X\) por \(f\) é o conjunto

	\[
		f[X]
		=
		\{
		f(x)
		\mid
		x\in X
		\}.
	\]

\end{defn}

\begin{obs}

	Alguns autores utilizam a notação

	\[
		f(X)
	\]

	em lugar de

	\[
		f[X].
	\]

\end{obs}

\begin{prop}
	\label{prop:matematica:funcoes:imagem-direta-e-inversa:imagem-direta-subconjunto}

	Para todo

	\[
		X\subseteq A,
	\]

	vale

	\[
		f[X]
		\subseteq
		B.
	\]

\end{prop}

\begin{proof}

	Se

	\[
		y\in f[X],
	\]

	então existe

	\[
		x\in X
	\]

	tal que

	\[
		y=f(x).
	\]

	Como

	\[
		f:A\to B,
	\]

	segue que

	\[
		y\in B.
	\]

\end{proof}

\begin{prop}
	\label{prop:matematica:funcoes:imagem-direta-e-inversa:monotonicidade-direta}

	Se

	\[
		X\subseteq Y\subseteq A,
	\]

	então

	\[
		f[X]
		\subseteq
		f[Y].
	\]

\end{prop}

\begin{proof}

	Seja

	\[
		z\in f[X].
	\]

	Então existe

	\[
		x\in X
	\]

	tal que

	\[
		z=f(x).
	\]

	Como

	\[
		X\subseteq Y,
	\]

	temos

	\[
		x\in Y.
	\]

	Logo

	\[
		z\in f[Y].
	\]

\end{proof}

\begin{prop}
	\label{prop:matematica:funcoes:imagem-direta-e-inversa:uniao-direta}

	Para quaisquer

	\[
		X,Y\subseteq A,
	\]

	vale

	\[
		f[X\cup Y]
		=
		f[X]
		\cup
		f[Y].
	\]

\end{prop}

\begin{proof}

	Segue diretamente das definições.

\end{proof}

\begin{prop}
	\label{prop:matematica:funcoes:imagem-direta-e-inversa:intersecao-direta}

	Para quaisquer

	\[
		X,Y\subseteq A,
	\]

	vale

	\[
		f[X\cap Y]
		\subseteq
		f[X]\cap f[Y].
	\]

\end{prop}

\begin{proof}

	Seja

	\[
		z\in f[X\cap Y].
	\]

	Então existe

	\[
		x\in X\cap Y
	\]

	tal que

	\[
		z=f(x).
	\]

	Como

	\[
		x\in X
	\]

	e

	\[
		x\in Y,
	\]

	segue que

	\[
		z\in f[X]
	\]

	e

	\[
		z\in f[Y].
	\]

	Portanto

	\[
		z\in f[X]\cap f[Y].
	\]

\end{proof}

\begin{obs}

	Em geral,

	\[
		f[X\cap Y]
		\neq
		f[X]\cap f[Y].
	\]

\end{obs}

\subsection{Imagem Inversa}

\begin{defn}[Imagem Inversa]
	\label{def:matematica:funcoes:imagem-direta-e-inversa:imagem-inversa}

	Seja

	\[
		f:A\to B
	\]

	e

	\[
		Y\subseteq B.
	\]

	A imagem inversa de \(Y\) por \(f\) é o conjunto

	\[
		f^{-1}[Y]
		=
		\{
		x\in A
		\mid
		f(x)\in Y
		\}.
	\]

\end{defn}

\begin{obs}

	A notação

	\[
		f^{-1}[Y]
	\]

	não pressupõe que \(f\) possua função inversa.

\end{obs}

\begin{prop}
	\label{prop:matematica:funcoes:imagem-direta-e-inversa:imagem-inversa-subconjunto}

	Para todo

	\[
		Y\subseteq B,
	\]

	vale

	\[
		f^{-1}[Y]
		\subseteq
		A.
	\]

\end{prop}

\begin{proof}

	Segue imediatamente da definição.

\end{proof}

\begin{prop}
	\label{prop:matematica:funcoes:imagem-direta-e-inversa:monotonicidade-inversa}

	Se

	\[
		Y_1\subseteq Y_2,
	\]

	então

	\[
		f^{-1}[Y_1]
		\subseteq
		f^{-1}[Y_2].
	\]

\end{prop}

\begin{proof}

	Seja

	\[
		x\in f^{-1}[Y_1].
	\]

	Então

	\[
		f(x)\in Y_1.
	\]

	Como

	\[
		Y_1\subseteq Y_2,
	\]

	segue que

	\[
		f(x)\in Y_2.
	\]

	Portanto

	\[
		x\in f^{-1}[Y_2].
	\]

\end{proof}

\begin{prop}
	\label{prop:matematica:funcoes:imagem-direta-e-inversa:uniao-inversa}

	Para quaisquer

	\[
		Y_1,Y_2\subseteq B,
	\]

	vale

	\[
		f^{-1}[Y_1\cup Y_2]
		=
		f^{-1}[Y_1]
		\cup
		f^{-1}[Y_2].
	\]

\end{prop}

\begin{proof}

	Segue diretamente da definição.

\end{proof}

\begin{prop}
	\label{prop:matematica:funcoes:imagem-direta-e-inversa:intersecao-inversa}

	Para quaisquer

	\[
		Y_1,Y_2\subseteq B,
	\]

	vale

	\[
		f^{-1}[Y_1\cap Y_2]
		=
		f^{-1}[Y_1]
		\cap
		f^{-1}[Y_2].
	\]

\end{prop}

\begin{proof}

	Segue diretamente da definição.

\end{proof}

\begin{prop}[Complemento]
	\label{prop:matematica:funcoes:imagem-direta-e-inversa:complemento}

	Para todo

	\[
		Y\subseteq B,
	\]

	vale

	\[
		f^{-1}[B\setminus Y]
		=
		A\setminus f^{-1}[Y].
	\]

\end{prop}

\begin{proof}

	Segue da definição de imagem inversa.

\end{proof}

\subsection{Caracterizações}

\begin{prop}[Caracterização da Injetividade]
	\label{prop:matematica:funcoes:imagem-direta-e-inversa:caracterizacao-injetividade}

	Seja

	\[
		f:A\to B.
	\]

	Então

	\[
		f
		\text{ é injetiva}
	\]

	se, e somente se,

	\[
		f^{-1}[f[X]]
		=
		X
	\]

	para todo

	\[
		X\subseteq A.
	\]

\end{prop}

\begin{proof}

	(\(\Rightarrow\))

	Seja

	\[
		x\in f^{-1}[f[X]].
	\]

	Então existe

	\[
		y\in X
	\]

	tal que

	\[
		f(x)=f(y).
	\]

	Pela injetividade,

	\[
		x=y.
	\]

	Logo

	\[
		x\in X.
	\]

	Portanto

	\[
		f^{-1}[f[X]]
		\subseteq
		X.
	\]

	A inclusão oposta é sempre válida.

	Logo há igualdade.

	(\(\Leftarrow\))

	Aplicar a hipótese ao conjunto

	\[
		X=\{x\}.
	\]

	A demonstração segue imediatamente.

\end{proof}

\begin{prop}[Caracterização da Sobrejetividade]
	\label{prop:matematica:funcoes:imagem-direta-e-inversa:caracterizacao-sobrejetividade}

	Seja

	\[
		f:A\to B.
	\]

	Então

	\[
		f
		\text{ é sobrejetiva}
	\]

	se, e somente se,

	\[
		f[f^{-1}[Y]]
		=
		Y
	\]

	para todo

	\[
		Y\subseteq B.
	\]

\end{prop}

\begin{proof}

	(\(\Rightarrow\))

	Seja

	\[
		y\in Y.
	\]

	Pela sobrejetividade existe

	\[
		x\in A
	\]

	tal que

	\[
		f(x)=y.
	\]

	Então

	\[
		x\in f^{-1}[Y].
	\]

	Consequentemente

	\[
		y\in f[f^{-1}[Y]].
	\]

	Logo

	\[
		Y
		\subseteq
		f[f^{-1}[Y]].
	\]

	A inclusão oposta é sempre válida.

	(\(\Leftarrow\))

	Tomando

	\[
		Y=B,
	\]

	obtemos

	\[
		f[f^{-1}[B]]
		=
		B.
	\]

	Como

	\[
		f^{-1}[B]=A,
	\]

	segue que

	\[
		f[A]=B.
	\]

	Portanto \(f\) é sobrejetiva.

\end{proof}

\end{document}
